
\linksubsection{tmp-path}{tmp\_path}
A \textbf{fresh empty directory} for every pytest invocation and every test.

\begin{itemize}[itemsep=0pt]
    \item Store generated input files for tested code
    \item Store output files, e.g.\ measurement data, \\ screenshots, etc.
    \item Access data even after pytest run is done, but no need for manual cleanup (3 directories are kept)
\end{itemize}

\begin{minted}[escapeinside=||]{python}
   from pathlib import Path

   def test_temp(|\highlightbox{tmp\_path: Path}|):
       |\codeskip|
\end{minted}
\vspace{.25em}%
\dirtree{%
    .1 /tmp/ \textsf{\footnotesize{\hspace{1em} \textcolor{gray}{(system-wide temp directory)}}}.
    .2 pytest-of-freya/.
    .3 \highlightbox[highlightcolorold]{pytest-1/}\vspace{-1pt}.
    .4 \highlightbox[highlightcolorold]{test\_temp0/}.
    .5 \codeskip*.
    .3 \highlightbox{pytest-2/}\vspace{-1pt}.
    .4 \highlightbox{test\_temp0/}.
    .5 \codeskip*.
    .3 \codeskip*.
}

\vspace{1.5em}
\mono{tmp_path}:

Based on Python's \mono{pathlib.Path}, a more object-oriented and convenient \mono{os.path} alternative:

\begin{minted}[autogobble]{python}
        dir_path = pathlib.Path("output")
        dir_path.mkdir(exist_ok=True)
        file_path = dir_path / "output.txt"
        file_path.write_text("Hello world!")
      \end{minted}

\bluefaicon{book} \monot{docs.python.org/3/library/pathlib.html}
\vspace{.65em}

\mono[codebluebright]{tmpdir}:

\textcolor{gray}{%
    Based on \mono[codebluebright]{py.path.local} from the \mono[codebluebright]{py} library (\mono[codebluebright]{pylib}) instead,
    since \mono[codebluebright]{pathlib} only exists since Python 3.4. Should not be used in new code.
}
